\section{共形块与射影平坦联络：这个对应是什么级别}
\label{sec:blocks}
\label{sec:conformal}
\label{sec:ch5-proj}

\subsection{这一章要解决什么}

第 3 章把三维 Chern--Simons 理论沿曲面 $\Sigma$ 切开，问题化为在 $\Sigma\times\mathbb R^1$ 上做正则量子化：约束 (3.3) 生成规范变换，经典相空间是 $\Sigma$ 上平坦 $G$ 联络的模空间 $\cM$ \src{370}。选一个复结构 $J$ 之后，$\cM$ 获得新的全纯解释——稳定全纯 $G_{\mathbb C}$ 丛的模空间 $\cM_J$；再取行列式线丛 $\cL$ 的 $k$ 次幂，量子 Hilbert 空间就是全纯截面空间
\begin{equation}
\cH_\Sigma^{(J)} \;=\; H^0\bigl(\cM_J,\;\cL^{\otimes k}\bigr) \src{369}.
\label{eq:ch5-hilbert}
\end{equation}
这个构造从头到尾用了 $J$，而物理问题只依赖光滑曲面 $\Sigma$。于是有四个递进的问题，本章就按它们展开。

\begin{enumerate}[leftmargin=1.6em]
\item[(Q1)] $\cH_\Sigma^{(J)}$ 随 $J$ 怎样变化？原文的回答是全纯地变化，因而可以看成一个复结构模空间上的全纯向量丛 \src{370}。
\item[(Q2)] 不同 $J$ 的纤维能不能\textbf{典范地}等同？原文说，这等价于该向量丛上存在一条典范平坦联络 \src{370}。
\item[(Q3)] 若不能整体等同，残余是什么？原文说这条联络只有\textbf{射影}平坦性，其单值性给出纯粹拓扑的映射类群作用 \src{370--371}。
\item[(Q4)] 这样得到的对应 $\Sigma\mapsto\cH_\Sigma$ 与 Segal 的模函子是什么关系？原文说“恰好重现”，但证据只到对象级 \src{370,396--397}。
\end{enumerate}

结论先放在这里，后文逐条落实。\textbf{(Q1)} 是真的，且是本章唯一不需要附加解释的部分：纤维构成全纯向量丛。\textbf{(Q2)} 的答案是否定的：因为共形反常，这条联络只到射影，Witten 自己在括号里作了限定，而且限定语“only”不能省 \src{370}。\textbf{(Q3)} 的答案是：单值性不是障碍，反而是拓扑数据；真正挡在“典范 $\cH_\Sigma$”前面的是射影因子，它使典范的对象只能是射影空间 $\mathbb P\cH_\Sigma$。\textbf{(Q4)} 的答案是：对象级等同（两边都是同一个共形块空间），函子级只有一句 “It is obvious” \src{397}，范畴级没有。因此本讲的正确措辞是“\textbf{强于维数相同，弱于范畴等价}”；这句话的精确内容见 \S\ref{sec:ch5-table} 的对照表。

\subsection{原文说了什么：p.369--371 的几段话}
\label{sec:ch5-quotes}

\paragraph{(a) 全纯依赖。}原文先把“依赖 $J$”说成事实，再把它解释成丛：
\begin{quote}\small
“As $J$ varies, the [\dots] vary holomorphically with $J$, and thus we could interpret this object as a holomorphic vector bundle on the moduli space of complex Riemann surfaces.” \src{370}
\end{quote}
这句话里没有证明，也没有用到共形场论的特殊性：它本质上是说，$\cM_J$、$\cL$ 以及全纯截面空间 $H^0$ 都随 $J$ 全纯地变动。用现代语言，这是 Grauert 直接像定理的一个实例（全纯族上的上同调构成全纯向量丛）；本讲义把它算作 \ours{} 的标准补充，原文没有引用任何定理。

\paragraph{(b) “与 $J$ 无关”被翻译成“有平坦联络”。}这是全章的方法论核心：
\begin{quote}\small
“The assertion that the [\dots] are canonically independent of $J$, and depend only on $\Sigma$, is the assertion that the vector bundle on moduli space [\dots] has a canonical flat connection that permits one to identify the fibers.” \src{370}
\end{quote}
把“典范等同”翻译成“平坦联络”，这一步是可逆的：给定模空间上的一条平坦联络，沿路径平行移动就给出纤维之间的线性等同；反过来，一族典范等同若光滑地依赖路径，就定义一条平坦联络。原文没有讨论这里的“典范”要不要与路径的选取无关；这一点在第 \S\ref{sec:ch5-proj} 节由单值性接管。

\paragraph{(c) 限定语：只有射影平坦，因子在定向反转下变号。}紧接着的括号是本章最不能省略的一句：
\begin{quote}\small
“The particular flat bundles on moduli space that we get this way are those that Segal has described [16] in connection with conformal field theory; Segal also rigorously proved the flatness, which is explained somewhat heuristically by the physical argument sketched above. (Because of the conformal anomaly, this bundle has only a projectively flat connection, with the projective factor being canonically odd under reversal of orientation.)” \src{370}
\end{quote}
这里有三层信息，必须分开读。(i) 这类平坦丛\textbf{就是} Segal 在共形场论里描述过的那一类，并且 Segal 严格证明了平坦性；Witten 自己的物理论证只是“somewhat heuristically”的解释。(ii) 由于共形反常，丛上只有\textbf{射影}平坦联络；“only”是限定，不是修辞。(iii) 射影因子在定向反转下变号，这是它“典范”的一部分——$\Sigma$ 与 $-\Sigma$ 的射影因子互为负值。

\paragraph{(d) 模函子，以及“恰好重现”。}原文随后给出这个对应的名字与语气：
\begin{quote}\small
“Segal calls the association [\dots] a `modular functor', and has given an algebra-geometric description of the modular functors that arise in current algebra. In quantizing the Chern-Simons theory we have exactly reproduced this description!” \src{370}
\end{quote}
“exactly reproduced”是强断言，但它说的对象是什么，要回到同页前半句：共形块被定义为 Ward 恒等式解的空间，“the space of conformal blocks” \src{370}。所以这句话至少断言：Chern--Simons 量子化给出的 $\cH_\Sigma^{(J)}$ 与流代数共形块是同一个向量空间。它没有断言缝合（gluing）相容，也没有断言配分函数层面的一致。

\paragraph{(e) 可以局部忘掉 $J$。}p.371 把 (b)(c) 合起来讲成操作：
\begin{quote}\small
“the existence of a natural projectively flat connection on the moduli space of complex structures permits one locally to (projectively) identify the various [\dots] and forget about the complex structure.” \src{371}
\end{quote}
注意括号里的 “(projectively)” 出现在“等同”之前，而不是之后：能忘掉的是射影类，不是向量本身。

\paragraph{(f) 单值性是映射类群，不是障碍。}同页紧接着处理“绕一圈回来怎么办”：
\begin{quote}\small
“these monodromies just correspond to an action of the purely topological mapping class group, so that the formal properties of [\dots] are just like those of $H^1(\Sigma,\mathbb R)$.” \src{371}
\end{quote}
原文先用 $\Sigma\mapsto H^1(\Sigma,\mathbb R)$ 作类比：这个函子不依赖复结构，微分同胚诱导线性变换，于是给出映射类群的表示。Chern--Simons 的 $\cH_\Sigma$ 应当具有同样的形式性质。这句话是 \S\ref{sec:ch5-proj} 的出发点：单值性即使是整体的，也只是拓扑的，不构成“依赖 $J$”的残余。

\paragraph{(g) 一处需要并读的措辞。}p.366 的脚注说，路径积分最自然给出的其实只是“$\cH_\Sigma$ 是与 $\Sigma$ 典范关联的向量空间，且在 $\Sigma$ 定向反转时与它的对偶互换” \src{366}。这与 p.370 的“只有射影平坦联络”并不矛盾，但两句的强度不同：p.366 是纲领性的（“典范关联的向量空间”），p.370 是精确化（典范的只是射影空间，且定向反转给对偶）。本讲义采用 p.370 的精确版本；把 p.366 读成“存在典范 $\cH_\Sigma$”是本章第一个要避免的误读。

\subsection{原文没说的细节：G5.1--G5.4}
\label{sec:ch5-gaps}

上一节的六段话给出了结论和语气，但没有给出定义、公式和归属。这一节补四件事：射影平坦的定量含义（G5.1）、这条联络在后世文献里的名字（G5.2）、平坦性到底由谁证明（G5.3）、以及“模函子”这个说法的内容与本文证据的差距（G5.4）。

\paragraph{G5.1 射影平坦到什么程度，射影因子从哪里来。}先把定义写清楚。设 $\pi:\mathcal{E}\to B$ 是复流形 $B$ 上的全纯向量丛，$\nabla$ 是全纯联络，$R^\nabla\in\Omega^2(B,\operatorname{End}\mathcal{E})$ 是曲率。称 $\nabla$ \textbf{射影平坦}，若存在 $2$-形式 $\omega\in\Omega^2(B)$ 使得
\begin{equation}
R^\nabla \;=\; \omega\cdot\operatorname{id}_{\mathcal{E}}.
\label{eq:ch5-projflat}
\end{equation}
也就是说，曲率取值在中心里；等价地，射影化丛 $\mathbb P\mathcal{E}$ 上有平坦联络，平行移动只在“差一个标量”的意义下与路径无关。注意这个定义与“平坦”只差一个标量 $2$-形式，但差别是本质的：$\omega\neq0$ 时不存在整体的平坦化，$\mathcal{E}$ 的纤维之间没有典范的线性等同，只有射影空间 $\mathbb P\mathcal{E}$ 的纤维之间有。

教材里对这一点有精确的定理。上野在共形块层 $V^\dagger(\mathcal C)$ 上构造联络 $\nabla^\omega$，并证明：\textbf{定理 5.5} 该联络是射影平坦的；证明的关键是曲率 $R$ 与提升的选取无关，并且“the curvature form $R$ may be regarded as a holomorphic two-form, hence the connection is projectively flat” [上野 p.93]。\textbf{定理 5.6} 给出曲率 $R(\xi,\eta)$ 的留数公式 [上野 p.95]。\textbf{推论 5.7}：当底空间是带 $N$ 个标记点的 $\mathbb P^1$ 或带 $N$ 个标记点的椭圆曲线的族时，该联络实际上是\textbf{平坦}的 [上野 p.95]。这条推论给出了一个判别机制：射影平坦的“射影”部分由底空间上的全纯 $2$-形式承载。在最简单的情形里这个机制可以一眼验证：$g=0$ 且 $n\le4$、或 $g=1$ 时底空间维数不超过 $1$，而复曲线上的 $\Omega^2=0$，故曲率必为零；一般的 $n$ 情形要依赖定理 5.6 的留数公式。$g\ge2$ 时底空间维数 $3g-3\ge3$，全纯 $2$-形式存在，曲率可以非零，这正是原文说 “only projectively” 的情形。

这里有一个上野文本没有明说、但由结论反推的细节：仅有“曲率是全纯 $2$-形式”还不够，必须再说明它乘的是恒等变换，即取值在中心。\ours{} 把这一点补成：曲率的标量性来自共形块层的构造（联络沿全纯方向的作用只改变截面的标量因子），因此 \eqref{eq:ch5-projflat} 成立；上野的证明末尾用 “hence the connection is projectively flat” 一句带过，正是默认了标量性。这个补充不影响任何结论，但它是“射影平坦”三个字的实际内容，不能省。

射影因子的物理来源，Witten 只给了一个名字：共形反常。原文的括号说 “Because of the conformal anomaly, this bundle has only a projectively flat connection” \src{370}，没有给公式，也没有解释反常怎样变成曲率。教材与本文合起来可以还原出机制：共形反常使沿模空间移动时的配分函数不是函数而是线丛 $\cL$ 的截面（这正是 \eqref{eq:ch5-hilbert} 的来源）；$k=1$ 时 $c_1(\cL)$ 就是辛形式，一般 $k$ 时是它的 $k$ 倍 \src{369}。因此曲率随 $k$ 线性增长，与中心荷 $c=k\dim G/(k+h^\vee)$ 的大 $k$ 行为一致。这个一致性只是量级上的核对，不是等式；把曲率写成 $c$ 乘 Weyl--Petersson 形式的精确公式需要归一化约定，原文没有给，本讲义也不把它当作已证结论。

最后是符号。原文说射影因子 “canonically odd under reversal of orientation” \src{370}。在教材的图像里这是可检验的：复结构反转 $J\mapsto-J$ 把 Kähler 形式 $\omega$ 变成 $-\omega$（第 3 章的辛形式 $\omega=(k/4\pi)\int_\Sigma\Tr(\delta A\wedge\delta A)$ 在 $J\mapsto-J$ 下变号），而曲率是全纯 $2$-形式、与 $\omega$ 成比例，于是 $R\mapsto-R$，射影因子取倒数即变号。\ours{} 这个推导只用到了辛形式的定义与 \eqref{eq:ch5-projflat}，不依赖任何未证命题；它解释了原文那句“典范地奇”的确切含义：不是任意一个符号，而是定向反转强制出来的符号。

\paragraph{G5.2 这条联络叫什么：Hitchin 联络与 KZ 联络。}1989 年的原文从头到尾没有给这条联络命名，也没有写出它的公式；它只说这类平坦丛是 Segal 描述过的那一类，并把严格性归于 Segal [16] \src{370}。后世文献把这条联络分成两种情形命名。

无标记点、$g\ge2$ 的情形叫 \textbf{Hitchin 联络}：N.~Hitchin, “Flat connections and geometric quantization”, \emph{Comm. Math. Phys.} \textbf{131} (1990), 347--380。Hitchin 在 Teichmüller 空间上构造了一族与 Kähler 极化相容的联络，证明它射影平坦，并用它证明几何量子化在复结构形变下的自然性。上野的 \textbf{定理 5.32} 说得很直接：“The Hitchin connection and our connection coincide” under the isomorphism (5.50) [上野 p.129]；这里的 “our connection” 就是共形块层上的联络。这是把“Witten 的射影平坦联络”与“Hitchin 联络”对上号的最干净的教材锚点。

带标记点、$g=0$ 的情形叫 \textbf{KZ 联络}（Knizhnik--Zamolodchikov，1984，\cite{knizhnik1984}）。它不是抽象的平坦性定理，而是显式公式：在 $\operatorname{Conf}_n(\mathbb P^1)$ 上，共形块丛的联络是
\begin{equation}
\nabla \;=\; \dd \;-\; \sum_{i=1}^{n} L^{(i)}_{-1}\,\dd z_i ,
\label{eq:ch5-kzconn}
\end{equation}
水平截面满足的方程就是 KZ 方程；在 $\mathfrak{sl}_2$ 情形，限制到最高权向量上得到
\begin{equation}
\frac{\partial\Psi_0}{\partial z_i}
\;=\; \frac{1}{k+2}\sum_{j\neq i}\frac{\Omega^{(ij)}\Psi_0}{z_i-z_j},
\label{eq:ch5-kz}
\end{equation}
其中 $\Omega^{(ij)}$ 是第 $i,j$ 个张量因子上的 Casimir 组合；这是一条每个分量都带分母 $k+2$ 的显式微分方程 [河野 p.27--28，式 (1.22) 与命题 1.17]。一般单群的对应分母是 $k+h^\vee$（上野式 (6.20)，$K=g^*+l$，[上野 p.130]；这里 $g^*$ 是对偶 Coxeter 数，$l$ 是 level）。

必须强调：\textbf{Hitchin 与 KZ 都是后世的归类，不是 1989 年原文的内容}。原文只引 Segal，不引 Hitchin，也不引 Knizhnik--Zamolodchikov；KZ 方程在正文中一次也没有出现（“Knizhnik” 在本译本中只出现在译注里）。因此在研讨班上正确的说法是：“Witten 的这条联络就是后来所谓的 Hitchin 联络（无标记点情形），其带标记点的显式形式是 KZ 联络”，而不是“Witten 使用了 Hitchin 联络”。

\paragraph{G5.3 平坦性到底由谁证明。}三处文献证明的是三件不同的事，不能互相替代。

(i) \textbf{Segal [16]}：证明的对象是模函子/共形块丛，Witten 的原话是 “Segal also rigorously proved the flatness” \src{370}。这是原文唯一承认的严格证明；Witten 自己的物理论证被他自己称为 “somewhat heuristically” \src{370}。引文见文献表 \cite{segal1988}。

(ii) \textbf{KZ / 河野 \S1.5}：证明的是 $\operatorname{Conf}_n(\mathbb P^1)$ 上显式联络 \eqref{eq:ch5-kzconn} 的平坦性。证明只有两步：系数 $L^{(i)}_{-1}$ 不依赖 $z$，故 $\dd\omega=0$；不同张量因子上的 $L$ 对易，故
\[
\omega\wedge\omega=\sum_{i<j}\bigl[L^{(i)}_{-1},L^{(j)}_{-1}\bigr]\,\dd z_i\wedge \dd z_j=0 ,
\]
于是曲率 $R=-\dd\omega+\omega\wedge\omega=0$ [河野 p.27--28]。这是本章唯一一个可以当场验证的平坦性证明，第 \S\ref{sec:ch5-deriv} 节把它完整写出。

(iii) \textbf{上野 \S5.1}：证明的是一般带点曲线族上共形块层的\textbf{射影}平坦性（定理 5.5、5.6），并在 $g=0,1$ 时给出平坦性（推论 5.7）[上野 p.91--95]。这比前两者更一般，但结论更弱——“射影”而非“平坦”，恰好对应 Witten 的括号。

三者的关系可以这样记：Segal 给的是公理层面的存在性，KZ 给的是 $\mathbb P^1$ 上的显式公式，上野给的是一般底空间上的曲率公式。本文只做了第四件事：用物理的配分函数论证把三者串起来，并把严格性推给 Segal。

\paragraph{G5.4 “模函子”的公理与本文证据的差距。}原文用 “modular functor” 命名 $\Sigma\mapsto\cH_\Sigma$，并说 Segal 对 current algebra 的模函子给了代数几何描述 \src{370}；p.366 又把模函子描述为 “canonically associate a Hilbert space to a Riemann surface” \src{366}。把这两句展开成公理（标准版本，\ours{} 据 Segal 的综述整理，原文没有逐条列出），至少有六条：

\begin{enumerate}[leftmargin=1.6em,itemsep=2pt]
\item[(M1)] \textbf{对象}：每个定向闭曲面 $\Sigma$（可带标记点或边界参数化）对应一个有限维复向量空间 $\cH_\Sigma$；
\item[(M2)] \textbf{定向反转}：$\cH_{-\Sigma}\cong\cH_\Sigma^{*}$，且这个等同是典范的；
\item[(M3)] \textbf{缝合}：沿边界粘合两个曲面时，$\cH$ 按融合规则作张量积并收缩；
\item[(M4)] \textbf{映射类群}：微分同胚类群作用在 $\cH_\Sigma$ 上，且只依赖拓扑；
\item[(M5)] \textbf{射影平坦丛}：$\cH_\Sigma$ 随复结构构成带射影平坦联络的丛，单值性给出 (M4)；
\item[(M6)] \textbf{融合/张量}：融合系数满足 Verlinde 型公式。
\end{enumerate}

把本文逐条对照：(M1) 成立，(M5) 成立，都是 p.369--371 的明确断言；(M2) 只在射影级别成立（p.366 脚注与 p.370 括号并读，$-\Sigma$ 给出对偶，但射影因子变号使“典范”只能到 $\mathbb P\cH$）；(M4) 成立，但原文只说 “formal properties \dots just like those of $H^1(\Sigma,\mathbb R)$” \src{371}，没有给出表示的具体构造；(M3) 是最大的缺口——原文在 \S4、\S5 用切开/粘合作论证，但缝合相容性只有 p.397 的 “It is obvious” \src{397}；(M6) 只有例子（三点函数维数 $N_{ijk}$，p.373）与第 9 章的 Verlinde 公式，本文未证。

因此对象级等同是有证据的：两边都等于 $H^0(\cM_J,\cL^{\otimes k})$，并且都带同一条射影平坦丛 \src{369--370}。函子级只有缝合这一条硬缺口；范畴级（自然变换、张量结构、模张量范畴）完全没有触及。这就是本章开头那句话的精确含义。

\subsection{补出的推导：D5.1--D5.3}
\label{sec:ch5-deriv}

这一节把本章三个最容易含糊的断言拆成可检验的步骤：为什么典范的对象只能是射影空间（D5.1）、KZ 联络为什么平坦以及系数 $1/(k+2)$ 从哪来（D5.2）、曲率为什么恰是射影因子并在什么情形消失（D5.3）。每个推导都注明依据；凡属本讲义补出的步骤，用 \ours{} 标出。

\paragraph{D5.1 从“只有射影平坦”到“只有 $\mathbb P\cH_\Sigma$ 典范”。}把 §\ref{sec:ch5-quotes} 的 (b)、(c) 两句当作前提：
\begin{itemize}[leftmargin=1.6em,itemsep=1pt]
\item[(P1)] “$\cH_\Sigma$ 典范地不依赖 $J$” 等价于 “该丛上有典范平坦联络” \src{370}；
\item[(P2)] 该丛上只有射影平坦联络，射影因子在定向反转下变号 \src{370}。
\end{itemize}
断言：存在典范的 $\cH_\Sigma$ 与 (P2) 矛盾；典范的对象是射影空间 $\mathbb P\cH_\Sigma$，定向反转在它上面给出对偶。证明分四步。

(1) \emph{若典范等同存在，则联络平坦。} 设对每对 $(J,J')$ 有典范线性等同 $E_{J,J'}:\cH_J\to\cH_{J'}$，且与路径的同伦相容。令 $\nabla$ 为“沿路径求导再比较”的联络，则平行移动就是 $E$，故 $\nabla$ 平坦。这一步正是 (P1) 的逆命题，见 \src{370} 的 “permits one to identify the fibers”。

(2) \emph{射影平坦的平行移动只到标量。} 由 (P2)，典范联络的曲率满足 \eqref{eq:ch5-projflat}，且 $\omega$ 不恒为零（否则联络平坦，与 “only” 矛盾）。于是沿路径的平行移动依赖路径的标量因子：$E_{J,J'}$ 只在“差一个非零标量”的意义下良定义，即它给出的是 $\mathbb P\cH_J\to\mathbb P\cH_{J'}$ 的映射。

(3) \emph{单值性落在 $\operatorname{PGL}$，而这不构成新障碍。} 绕闭合路径一周，单值性 $\rho(\gamma)\in\operatorname{PGL}(\cH_\Sigma)$。原文 p.371 说这些单值性就是映射类群的作用 \src{371}；这完全相容：$\rho$ 非平凡并不违反 (P2)，(P2) 禁止的只是典范的\emph{线性}提升。反过来，若有人声称有典范线性等同，就得到一个平坦联络，与 (P2) 冲突。所以结论只能是 $\mathbb P\cH_\Sigma$ 典范。

(4) \emph{为什么不能靠“选一个提升”补救。} 把射影平坦联络提升为线性平坦联络，等价于选一个整体 $1$-形式 $\theta$ 使 $\dd\theta=\omega$，或等价地选丛的一个平坦化；两个这样的选择相差闭 $1$-形式，其积分给出标量。定向反转 $J\mapsto-J$ 把 $\omega$ 变成 $-\omega$（G5.1 已算），故提升的标量因子也变号；这就是 “projective factor being canonically odd under reversal of orientation” \src{370} 的确切内容，也是 p.366 脚注 “exchanged with its dual” 在射影级别的版本 \src{366}。\ours{} 这一步的可检验后果是：任何声称“有典范 $\cH_\Sigma$”的构造，都必须给出一个在定向反转下不变的提升；而 (P2) 的奇性禁止这样的提升。

\begin{keypoint}
“射影平坦”保证：射影空间纤维之间有典范等同；单值性是 $\operatorname{PGL}$ 中的映射类群表示；定向反转时射影因子变号。它\emph{不}保证：向量纤维之间有典范等同；存在整体平坦化；单值性平凡。把“单值性非平凡”当作“不能忘掉 $J$”的理由是错的——原文 p.371 恰恰说单值性只是拓扑作用。
\end{keypoint}

\paragraph{D5.2 KZ 联络：显式公式与平坦性。}下面的推导把 \eqref{eq:ch5-kzconn}--\eqref{eq:ch5-kz} 完整走一遍，并核对系数的来源。

\begin{derivation}[KZ 联络为什么平坦，$k+2$ 从哪来]
\emph{设定.} $g=0$，$n$ 个标记点 $z=(z_1,\dots,z_n)\in\operatorname{Conf}_n(\mathbb P^1)$；纤维是共形块，即多重线性映射在不变量子代数上不变的空间 $\cH_z=\Hom_{\mathfrak g(z)}(\otimes_i H_{\lambda_i},\mathbb C)$ [河野 p.19, p.25]。

\emph{第一步（联络）.} 定义 $\nabla=\dd-\omega$，其中
\[
\omega=\sum_{i=1}^{n}L^{(i)}_{-1}\,\dd z_i ,
\]
$L^{(i)}_{-1}$ 是第 $i$ 个张量因子上仿射代数的降算子。水平截面满足 $\partial_i\Psi=L^{(i)}_{-1}\Psi$，这就是河野书中的 (1.22) [河野 p.27]。

\emph{第二步（$\dd\omega=0$）.} $L^{(i)}_{-1}$ 是常算子，不依赖 $z$，故
\[
\dd\omega=\sum_i \dd\bigl(L^{(i)}_{-1}\bigr)\wedge\dd z_i=0 .
\]

\emph{第三步（$\omega\wedge\omega=0$）.} 按定义展开：
\[
\omega\wedge\omega=\sum_{i,j}L^{(i)}_{-1}L^{(j)}_{-1}\,\dd z_i\wedge\dd z_j
=\sum_{i<j}\bigl[L^{(i)}_{-1},L^{(j)}_{-1}\bigr]\,\dd z_i\wedge\dd z_j=0,
\]
因为 $i\neq j$ 时两个算子作用在不同张量因子上，对易。

\emph{第四步（曲率）.} 对 $\nabla=\dd-\omega$ 有 $R=-\dd\omega+\omega\wedge\omega$，由前两步 $R=0$。故 $\nabla$ 平坦 [河野 p.27--28]。

\emph{第五步（化为 KZ 方程）.} 把水平方程限制到最高权向量，用 Sugawara 构造写出 $L^{(i)}_{-1}$ 的显式作用：
\begin{equation}
\frac{\partial\Psi_0}{\partial z_i}
=\frac{1}{k+2}\sum_{j\neq i}\frac{\Omega^{(ij)}\Psi_0}{z_i-z_j},
\qquad \Omega^{(ij)}=\sum_a T^a_{(i)}T^a_{(j)},
\label{eq:ch5-kz2}
\end{equation}
即河野的命题 1.17 [河野 p.28]。

\emph{第六步（两点检验）.} 取 $n=2$、$G=\operatorname{SU}(2)$、表示 $j$。共形块是单态，一维；由 $\vec J_1+\vec J_2=0$ 与 $C_2=j(j+1)$ 得 $\Omega^{(12)}=-C_2(j)=-j(j+1)$（通常归一化）。于是 \eqref{eq:ch5-kz2} 变成
\[
\frac{\partial\Psi}{\partial z}=-\frac{j(j+1)}{k+2}\cdot\frac{\Psi}{z-w}
\;\Longrightarrow\;
\Psi\propto(z-w)^{\,j(j+1)/(k+2)}=(z-w)^{h_j},
\]
恰是共形权 $h_j=j(j+1)/(k+2)$ 的幂律。\guess{} 不同教材对 $\Omega$ 的整体符号约定不同（河野取使上式右端为正的约定），差一个符号不影响指数；指数等于 $h_j$ 这一点与约定无关。

\emph{第七步（一般单群）.} 把 $k+2$ 换成 $k+h^\vee$：$h^\vee$ 是对偶 Coxeter 数，$\kappa=k+h^\vee$ 是 KZ 的分母。河野在批注中给出对应 $k+2\leftrightarrow k+h^\vee$（[河野 p.107]），上野的式 (6.20) 写作 $K=g^*+l$（$g^*$ 为对偶 Coxeter 数、$l$ 为 level）[上野 p.130]。两点检验同样给出 $h_\lambda=(\lambda,\lambda+2\rho)/(2(k+h^\vee))$。

\emph{适用范围.} 本推导只覆盖 $g=0$；$g\ge1$ 时 $\operatorname{Conf}$ 上的曲率不再自动为零，推导止于“射影平坦”，见 D5.3。
\end{derivation}

\paragraph{D5.3 曲率为什么是射影因子，何时消失。} (a) 上野定理 5.5 的证明结构是：曲率 $R(\xi,\eta)$ 与向量场的提升无关（否则联络不良定）；沿全纯方向的协变导数与 $\bar\partial$ 相容，故 $R$ 可以看作底空间上的全纯 $2$-形式；作为 $\operatorname{End}$ 值形式，它只能是这个 $2$-形式的标量倍数 [上野 p.93]。\ours{} 最后这句“标量倍数”是射影平坦定义 \eqref{eq:ch5-projflat} 的直接翻译，上野文本用 “hence the connection is projectively flat” 一句带过。

(b) 何时消失：$g=0$ 且 $n\le4$、或 $g=1$ 时底空间维数不超过 $1$，复曲线上 $\Omega^2=0$，故 $R=0$，联络平坦（推论 5.7，[上野 p.95]）；一般的 $g=0,n\ge5$ 情形由定理 5.6 的留数公式给出消失，而不是仅由类型理由。$g\ge2$ 时底空间维数 $3g-3\ge3$，全纯 $2$-形式存在，曲率可以非零；原文对一般 $\Sigma$ 断言 “only a projectively flat connection” \src{370}，指的就是这种情形。

(c) 定向反转的符号：$J\mapsto-J$ 把辛形式 $\omega$ 变成 $-\omega$（第 3 章的 $\omega=(k/4\pi)\int_\Sigma\Tr\delta A\wedge\delta A$ 中 $J$ 只进入 $*$ 与复结构，反转给出整体负号）。若曲率与 $\omega$ 成比例，$R\mapsto-R$，射影因子取倒数即变号；与 p.370 括号的 “odd” 一致，也与 p.366 脚注“定向反转给对偶”一致。

(d) 反常的角色：把比例常数写成中心荷 $c=k\dim G/(k+h^\vee)$ 是 \ours{} 的重建，没有原文支持到“等式”的程度；能核对的只是量级：$k=1$ 时 $c_1(\cL)$ 就是辛形式，一般 $k$ 时为 $k$ 倍 \src{369}，故曲率随 $k$ 线性增长，与 $c$ 的大 $k$ 行为一致。原文只写 “Because of the conformal anomaly” \src{370}，没有公式；本讲义不把精确常数当作已证结论。

\subsection{与其他来源对照}
\label{sec:ch5-table}

\begin{table}[htbp]
\centering\small
\begin{tabular}{@{}p{2.2cm}p{5.4cm}p{2.2cm}p{4.2cm}@{}}
\toprule
来源 & 断言 / 对象 & 级别 & 锚点\\
\midrule
Witten 1989 & 对象级等同；随 $J$ 全纯；只有射影平坦；单值性 = 映射类群；引 Segal & 物理 + 引证 & \src{369--371}\\
Segal [16] & 模函子 / 流代数共形块；严格证明平坦性 & 定理 & p.370 引；\cite{segal1988}\\
河野 \S1.5 & KZ 联络 $\nabla=\dd-\sum_i L^{(i)}_{-1}\dd z_i$；$\dd\omega=0$、$\omega\wedge\omega=0$ 故平坦 & 定理 + 证明 & [河野 p.25--28]\\
上野 \S5.1 & 共形块层的射影平坦联络；曲率 = 全纯 $2$-形式；$g=0,1$ 时平坦 & 定理 5.5/5.6、推论 5.7 & [上野 p.91--95]\\
上野 \S5.6 & Hitchin 联络与共形块联络一致 & 定理 5.32 & [上野 p.129]\\
Hitchin 1990 & 无标记点、Teichmüller 空间上的射影平坦联络 & 定理（后世） & \emph{CMP} \textbf{131} (1990) 347--380\\
李思 5.4 & Kähler 极化 $\Rightarrow$ P-极化截面 = 全纯截面 & 教材框架 & [李思 p.304--309]\\
本文 \S5 & “can be interpreted as”“It is obvious” & 解释 / 纲领 & \src{396--397}\\
\bottomrule
\end{tabular}
\caption{共形块与射影平坦联络：各来源的断言与证据级别。前四行是本章的主证据链；第五、六行是后世的命名与统一；第七行给出几何量子化的一般框架；最后一行是本文自身的最强措辞。}
\label{tab:ch5-compare}
\end{table}

读这张表要注意两点。第一，\textbf{证明级别不同}：Segal 的平坦性、河野的 KZ 平坦性、上野的射影平坦性是三条独立的定理，证的对象分别是模函子、$\operatorname{Conf}_n(\mathbb P^1)$ 上的显式丛、一般带点曲线族上的共形块层；本文没有证明其中任何一条，只做了物理的串联。第二，\textbf{“平坦”与“射影平坦”在不同底空间上互换}：$g=0,1$ 时上野的联络平坦（推论 5.7），与河野的 KZ 平坦性一致；$g\ge2$ 时只有射影平坦，这正是 Witten 括号所指的情形。把这两点混在一起，就会得出“Witten 证明了联络平坦”或“KZ 与 CS 是同一个理论”这类过强的结论。

\subsection{陷阱与约定}
\label{sec:ch5-traps}

\begin{enumerate}[leftmargin=1.7em,itemsep=3pt]
\item \textbf{“与 $J$ 无关”不等于“有典范 $\cH_\Sigma$”。} p.366 脚注的 “canonically associated” 是纲领性措辞，p.370 的括号把它精确化为“只有射影平坦”。典范的对象是 $\mathbb P\cH_\Sigma$，不是 $\cH_\Sigma$；定向反转给出的是射影对偶。
\item \textbf{“射影平坦”不等于“平坦”。} $g=0$（$n\le4$）与 $g=1$ 时是平坦的（上野推论 5.7），$g\ge2$ 时只有射影平坦。“only a projectively flat connection” 里的 “only” 是限定语，不能改写成 “a flat connection”。
\item \textbf{射影因子不是单值性。} 单值性是纯粹拓扑的映射类群作用 \src{371}，它不构成“忘不掉 $J$”的残余；射影因子的来源是共形反常 \src{370}。把两者混为一谈，就会误以为“单值性非平凡所以只能射影”。
\item \textbf{“Segal 证明了平坦性”证的是什么。} 证的是模函子 / 流代数共形块丛的平坦性，不是 Chern--Simons 路径积分的严格数学构造；Witten 自己说物理论证 “somewhat heuristically” \src{370}。同句把 “flatness” 与 “projectively flat” 并列，措辞有张力；本讲义把它读作：Segal 的定理已把射影因子作为模函子结构的一部分，因此“平坦”是在射影意义下说的。
\item \textbf{Hitchin 与 KZ 是后世的归类。} 1989 年原文没有这两个名字，也没有 KZ 方程；“Knizhnik” 在本译本中只出现在译注里。正确说法是“后来被称为 Hitchin 联络 / 其带标记点的显式形式是 KZ 联络”。
\item \textbf{不要混淆两个分母。} KZ 的分母 $\kappa=k+h^\vee$ 是共形权与 KZ 系数的组合参数（第 4 章已从 $\eta$ 因子得到同一有效能级）；第 6 章的可积性条件 $2j\le k$ 用的却是裸 level $k$。两者语境不同，不能互相代入。
\item \textbf{不要写“范畴等价”。} 本文最强的是对象级等同（两边都是 $H^0(\cM_J,\cL^{\otimes k})$ 并带同一射影平坦丛）；函子级缺缝合，范畴级没有。正确措辞是“强于维数相同，弱于范畴等价”。
\item \textbf{定向反转的等同只在射影级。} $\cH_{-\Sigma}\cong\cH_\Sigma^{*}$ 不能写成线性典范等同；射影因子在反转下变号，正是这个“只能到射影”的体现。
\item \textbf{模函子公理不是原文。} 本章 \S\ref{sec:ch5-gaps} 列的 (M1)--(M6) 是标准版本，据 Segal 的综述整理（\ours）；原文只给了名字与一句描述。引用公理时要说明来源，不能当作 Witten 的原文。
\end{enumerate}

\subsection{练习}
\label{sec:ch5-exercises}

\begin{enumerate}[leftmargin=1.8em,itemsep=3pt]
\item \textbf{（定位）} 在 p.370 找出三句话：一句是断言（“$J$ 无关”等价于“平坦联络”），一句是引证（Segal 证明了平坦性），一句是限定（只有射影平坦）。说明为什么删掉限定句会改变整章的结论。\emph{提示}：对照 p.366 脚注与 p.371 的 “(projectively) identify”。
\item \textbf{（定义）} 设 $\nabla$ 是 $B$ 上向量丛的联络，曲率 $R=\omega\cdot\operatorname{id}$，$\omega\neq0$。证明：在 $B$ 的单连通覆盖上存在平坦化 $\nabla+\theta$（即 $\dd\theta=\omega$），但 $\theta$ 没有典范选择；由此说明为什么平行移动只给出 $\mathbb P\mathcal{E}$ 的典范等同。
\item \textbf{（KZ）} 按 D5.2 的步骤，对 $G=\operatorname{SU}(2)$、$n=2$ 解 KZ 方程，验证解是指数 $h_j=j(j+1)/(k+2)$ 的幂律。再说明：证明平坦性时哪一步用到了“不同张量因子上的算子对易”，如果 $L^{(i)}_{-1}$ 依赖 $z_i$ 会发生什么。
\item \textbf{（判据）} 用上野推论 5.7 判断下列四种情形哪些是平坦、哪些只射影平坦，并给理由：$g=0,n=3$；$g=0,n=5$；$g=1,n=0$；$g=2,n=0$。\emph{提示}：看底空间维数与全纯 $2$-形式，并注意 $n\ge5$ 时要引用定理 5.6。
\item \textbf{（反常）} 计算 $\operatorname{SU}(2)$、$k=1$ 与 $\operatorname{SU}(3)$、$k=2$ 的中心荷 $c=k\dim G/(k+h^\vee)$，并与 p.369 的 $c_1(\cL)=\omega$（$k=1$）核对量级。说明为什么这个计算只能支持“曲率随 $k$ 线性”的量级结论，不能当作精确公式。
\item \textbf{（公理）} 对 \S\ref{sec:ch5-gaps} 的 (M1)--(M6) 逐条标注本文的证据等级（有断言 / 只有 “obvious” / 没有），并指出哪一条是函子级对应的最大缺口。\emph{提示}：p.397 的 “It is obvious” 对应哪一条？
\item \textbf{（辨析）} 有人论证：“单值性非平凡，所以不能忘掉 $J$，所以只有射影平坦。”用 p.371 的原话指出这个论证错在哪一步，并给出正确的关系：单值性、射影因子、反常三者各起什么作用。
\item \textbf{（选做）} 把 p.366 脚注与 p.370 括号放在一起，写出 $\cH_{-\Sigma}$ 与 $\cH_\Sigma^{*}$ 的精确关系（到射影级），并说明射影因子变号在这个关系里如何体现。
\end{enumerate}

\subsection{本章边界}
\label{sec:ch5-boundary}

\begin{boundary}
\textbf{本章确立（有原文或教材锚点）}：(i) $\cH_\Sigma^{(J)}=H^0(\cM_J,\cL^{\otimes k})$ 随 $J$ 全纯地构成丛 \src{369--370}；(ii) “典范不依赖 $J$” 等价于“典范平坦联络” \src{370}；(iii) 该联络只有射影平坦性，射影因子在定向反转下变号 \src{370}；(iv) 单值性是纯粹拓扑的映射类群作用 \src{371}；(v) 这类丛就是 Segal 描述的那一类，平坦性由 Segal 严格证明 \src{370}；(vi) KZ 联络的显式形式与平坦性 [河野 p.25--28]，一般单群的分母 $k+h^\vee$ [上野 p.130]；(vii) 共形块联络的射影平坦性与曲率公式 [上野 p.91--95]，Hitchin 联络与之相同 [上野 p.129]。

\textbf{本章补出（\ours）}：射影平坦定义 \eqref{eq:ch5-projflat} 与“只有 $\mathbb P\cH$ 典范”的四步论证（D5.1）；射影因子变号的符号推导（G5.1、D5.3(c)）；KZ 两点检验（D5.2 第六步）；中心荷与曲率的量级核对（D5.3(d)）；模函子公理 (M1)--(M6) 与本文证据的逐条对照（G5.4）。

\textbf{未解决 / 未触及}：缝合相容性只有 p.397 的 “It is obvious”，本文没有证明；函子级对应缺一条硬证据；范畴级（自然变换、张量结构、模张量范畴）完全没有；射影因子的精确常数与半形式修正的符号本文未给；Segal [16] 的严格证明本身没有被复核，只按原文引用。

\textbf{待核}：p.370 同一句里 “flatness” 与 “only a projectively flat connection” 的并置，究竟是 Segal 定理已含射影因子，还是两者针对相差一个线丛的不同的丛——本讲义取前者（第 4 条陷阱），但未在原始文献中逐字核对。
\end{boundary}



